\section{Methods}

\subsection{Characterizing rice-atlas}
We inspected \texttt{dr-richard-barker/rice-atlas} via the GitHub API and \texttt{gh} CLI
(repository metadata, full file tree, Pages configuration, and commit history) rather than
assuming its content from its name or description. This established that it is a fork
(itself forked from a human-anatomy viewer template) with no commits from the account that
owns it, no \texttt{data/} directory, no cited primary dataset, and a license the GitHub
license detector could not resolve to a standard SPDX identifier.

\subsection{Provenance discipline}
Every external fact used in this repository is recorded in \texttt{data/README.md} with
its source, verification method, and license status before being used anywhere else. Three
literature DOIs (\citealp{Klepikova2016,Shahan2022,Vijayan2021}) were each confirmed
against the CrossRef API (title, full author list, year, volume, pages) rather than
trusted from a search snippet. Four NASA OSDR study accessions were confirmed live against
the OSDR metadata API before use; two (OSD-120, OSD-314) had already-fetched raw count
matrices available (reused, with permission implicit in common ownership, from the sibling
repository \texttt{arabidopsis-drem-osdr}, itself sourced from the same public, CC0 OSDR
accessions) and were used quantitatively; two (OSD-38, OSD-522) are cited as confirmed,
relevant, real studies but are not yet quantitatively integrated, and are documented as
pending rather than fabricated.

\subsection{Spaceflight response quantification}
For each of the two studies with available raw counts, samples were split into their real
experimental groups by the labels already present in the count matrix's own sample names
(OSD-120: \texttt{GC} ground control vs.\ \texttt{FLT} flight, Day 13 root tissue; OSD-314:
\texttt{1g} vs.\ \texttt{0g}, whole seedling). Raw counts were converted to counts-per-million
(CPM) per sample to correct for sequencing-depth differences, averaged within each group,
and reported as $\log_2\!\big((\text{mean CPM}_{\text{treatment}}+1)/(\text{mean
CPM}_{\text{control}}+1)\big)$. This is deliberately a simple mean-ratio summary: no
DESeq2/edgeR model was fit, no $p$-values or multiple-testing correction were computed, and
no covariates (genotype, light treatment) were controlled for. The implementation is
\texttt{scripts/01\_compute\_spaceflight\_response.py}; its output feeds
\texttt{scripts/02\_export\_viewer\_data.py}, which derives the compact per-organ summary
shown in the viewer, so no number reaching the viewer or this manuscript is hand-typed.

\subsection{Structural model}
Organ geometry for the viewer (root system, rosette leaf, inflorescence axis, flower,
silique) is generated procedurally in the client (\texttt{app/src/geometry.ts}) using
Three.js primitive shapes, not sculpted or scanned. Root branching topology and the
rosette's leaf-arrangement angle follow parameters attributed, organ by organ, to the
literature sources in \texttt{app/src/organs.ts}; where a parameter is illustrative rather
than measured, that organ's on-screen description says so explicitly, matching the
disclosure standard rice-atlas itself already sets for its own geometry.

\subsection{Evidence gate}
Before this repository's first push, and again before any push touching
publication-relevant content, we ran this portfolio's own pre-publication check (a
blocklist sweep for previously-fabricated strings, CrossRef verification of every DOI,
confirmation that every number traces to a producing script with no synthetic-data
fallback, and confirmation that every cited file path and accession resolves in the actual
repository). Only a check that was actually performed is recorded as an attestation; an
attestation is never written to unblock a push it did not actually clear.
